\documentclass[12pt,a4paper]{article}
\usepackage[utf8]{inputenc}
\usepackage{amsmath,amssymb,amsthm}
\usepackage{physics}
\usepackage{geometry}
\usepackage{tikz}
\usepackage{pgfplots}
\usepackage{hyperref}
\usepackage{enumitem}
\usepackage{xcolor}
\usepackage{fancyhdr}
\usepackage{tcolorbox}
\usepackage{multicol}

\geometry{margin=1in}
\pagestyle{fancy}
\fancyhf{}
\fancyhead[L]{\textbf{GATE Mathematical Physics Study Guide}}
\fancyhead[R]{\thepage}
\fancyfoot[C]{2026 Preparation}

\hypersetup{
    colorlinks=true,
    linkcolor=blue,
    filecolor=magenta,      
    urlcolor=cyan,
}

\newtcolorbox{importantbox}{
    colback=red!5!white,
    colframe=red!75!black,
    title=\textbf{Important},
    fonttitle=\bfseries
}

\newtcolorbox{formulabox}{
    colback=blue!5!white,
    colframe=blue!75!black,
    title=\textbf{Key Formula},
    fonttitle=\bfseries
}

\newtcolorbox{examplebox}{
    colback=green!5!white,
    colframe=green!75!black,
    title=\textbf{Example},
    fonttitle=\bfseries
}

\title{\textbf{GATE 2026: Mathematical Physics Study Guide}\\
\large Comprehensive Coverage - Leave No Stone Unturned}
\author{Preparation Material}
\date{2026}

\begin{document}

\maketitle
\tableofcontents
\newpage

\section{Introduction}
Mathematical Physics is the foundation of all physics. Mastery of these topics is essential for success in GATE. This guide provides comprehensive coverage of:
\begin{itemize}
    \item Vector Calculus
    \item Complex Analysis
    \item Differential Equations
    \item Linear Algebra
\end{itemize}

\newpage

% ============================================
% VECTOR CALCULUS
% ============================================
\section{Vector Calculus}
Vector calculus is fundamental for understanding fields, flux, and circulation in physics.

\subsection{Vector Operations and Notation}
\subsubsection{Basic Definitions}
\begin{itemize}
    \item \textbf{Vector field}: $\vec{F}(x,y,z) = F_x\hat{i} + F_y\hat{j} + F_z\hat{k}$
    \item \textbf{Position vector}: $\vec{r} = x\hat{i} + y\hat{j} + z\hat{k}$
    \item \textbf{Magnitude}: $|\vec{F}| = \sqrt{F_x^2 + F_y^2 + F_z^2}$
\end{itemize}

\subsubsection{Vector Products}
\begin{formulabox}
\textbf{Scalar (Dot) Product:}
\begin{align}
\vec{A} \cdot \vec{B} &= A_x B_x + A_y B_y + A_z B_z \\
&= |\vec{A}||\vec{B}|\cos\theta
\end{align}

\textbf{Vector (Cross) Product:}
\begin{align}
\vec{A} \times \vec{B} &= \begin{vmatrix}
\hat{i} & \hat{j} & \hat{k} \\
A_x & A_y & A_z \\
B_x & B_y & B_z
\end{vmatrix} \\
&= |\vec{A}||\vec{B}|\sin\theta \hat{n}
\end{align}
\end{formulabox}

\subsection{Gradient (Grad)}
The gradient of a scalar field $\phi(x,y,z)$ is a vector field that points in the direction of the greatest rate of increase.

\begin{formulabox}
\textbf{Gradient in Cartesian coordinates:}
\begin{equation}
\grad \phi = \pdv{\phi}{x}\hat{i} + \pdv{\phi}{y}\hat{j} + \pdv{\phi}{z}\hat{k}
\end{equation}

\textbf{Alternative notation:}
\begin{equation}
\grad \phi = \left(\hat{i}\pdv{}{x} + \hat{j}\pdv{}{y} + \hat{k}\pdv{}{z}\right)\phi
\end{equation}
\end{formulabox}

\subsubsection{Properties of Gradient}
\begin{enumerate}
    \item \textbf{Linearity}: $\grad(\phi + \psi) = \grad\phi + \grad\psi$
    \item \textbf{Product rule}: $\grad(\phi\psi) = \phi\grad\psi + \psi\grad\phi$
    \item \textbf{Chain rule}: $\grad f(\phi) = f'(\phi)\grad\phi$
    \item \textbf{Directional derivative}: $D_{\hat{u}}\phi = \hat{u} \cdot \grad\phi$
\end{enumerate}

\subsubsection{Gradient in Other Coordinate Systems}
\begin{formulabox}
\textbf{Cylindrical coordinates} $(r, \theta, z)$:
\begin{equation}
\grad \phi = \pdv{\phi}{r}\hat{r} + \frac{1}{r}\pdv{\phi}{\theta}\hat{\theta} + \pdv{\phi}{z}\hat{z}
\end{equation}

\textbf{Spherical coordinates} $(r, \theta, \phi)$:
\begin{equation}
\grad \phi = \pdv{\phi}{r}\hat{r} + \frac{1}{r}\pdv{\phi}{\theta}\hat{\theta} + \frac{1}{r\sin\theta}\pdv{\phi}{\phi}\hat{\phi}
\end{equation}
\end{formulabox}

\begin{examplebox}
\textbf{Example:} Find the gradient of $\phi(x,y,z) = x^2y + yz^2$.

\textbf{Solution:}
\begin{align}
\grad\phi &= \pdv{}{x}(x^2y + yz^2)\hat{i} + \pdv{}{y}(x^2y + yz^2)\hat{j} + \pdv{}{z}(x^2y + yz^2)\hat{k} \\
&= 2xy\hat{i} + (x^2 + z^2)\hat{j} + 2yz\hat{k}
\end{align}
\end{examplebox}

\subsection{Divergence (Div)}
The divergence of a vector field measures the "outflow" of the field from a point.

\begin{formulabox}
\textbf{Divergence in Cartesian coordinates:}
\begin{equation}
\div \vec{F} = \pdv{F_x}{x} + \pdv{F_y}{y} + \pdv{F_z}{z}
\end{equation}

\textbf{Using del operator:}
\begin{equation}
\div \vec{F} = \grad \cdot \vec{F} = \left(\hat{i}\pdv{}{x} + \hat{j}\pdv{}{y} + \hat{k}\pdv{}{z}\right) \cdot \vec{F}
\end{equation}
\end{formulabox}

\subsubsection{Physical Interpretation}
\begin{itemize}
    \item $\div \vec{F} > 0$: Source point (field is diverging)
    \item $\div \vec{F} < 0$: Sink point (field is converging)
    \item $\div \vec{F} = 0$: Solenoidal field (no sources/sinks)
\end{itemize}

\subsubsection{Properties of Divergence}
\begin{enumerate}
    \item \textbf{Linearity}: $\div(\vec{F} + \vec{G}) = \div\vec{F} + \div\vec{G}$
    \item \textbf{Product rule}: $\div(\phi\vec{F}) = \phi\div\vec{F} + \vec{F} \cdot \grad\phi$
    \item \textbf{Divergence of gradient}: $\div(\grad\phi) = \laplacian\phi$ (Laplacian)
\end{enumerate}

\subsubsection{Divergence in Other Coordinate Systems}
\begin{formulabox}
\textbf{Cylindrical coordinates:}
\begin{equation}
\div \vec{F} = \frac{1}{r}\pdv{}{r}(rF_r) + \frac{1}{r}\pdv{F_\theta}{\theta} + \pdv{F_z}{z}
\end{equation}

\textbf{Spherical coordinates:}
\begin{equation}
\div \vec{F} = \frac{1}{r^2}\pdv{}{r}(r^2F_r) + \frac{1}{r\sin\theta}\pdv{}{(\sin\theta F_\theta)} + \frac{1}{r\sin\theta}\pdv{F_\phi}{\phi}
\end{equation}
\end{formulabox}

\begin{examplebox}
\textbf{Example:} Find the divergence of $\vec{F} = x^2\hat{i} + xy\hat{j} + xz\hat{k}$.

\textbf{Solution:}
\begin{align}
\div\vec{F} &= \pdv{}{x}(x^2) + \pdv{}{y}(xy) + \pdv{}{z}(xz) \\
&= 2x + x + x = 4x
\end{align}
\end{examplebox}

\subsection{Curl}
The curl of a vector field measures the rotation or "circulation" of the field.

\begin{formulabox}
\textbf{Curl in Cartesian coordinates:}
\begin{equation}
\curl \vec{F} = \begin{vmatrix}
\hat{i} & \hat{j} & \hat{k} \\
\pdv{}{x} & \pdv{}{y} & \pdv{}{z} \\
F_x & F_y & F_z
\end{vmatrix}
\end{equation}

\textbf{Expanded form:}
\begin{equation}
\curl \vec{F} = \left(\pdv{F_z}{y} - \pdv{F_y}{z}\right)\hat{i} + \left(\pdv{F_x}{z} - \pdv{F_z}{x}\right)\hat{j} + \left(\pdv{F_y}{x} - \pdv{F_x}{y}\right)\hat{k}
\end{equation}
\end{formulabox}

\subsubsection{Physical Interpretation}
\begin{itemize}
    \item $\curl \vec{F} = \vec{0}$: Irrotational field (conservative)
    \item $\curl \vec{F} \neq \vec{0}$: Rotational field (non-conservative)
    \item Magnitude of curl: measures the strength of rotation
\end{itemize}

\subsubsection{Properties of Curl}
\begin{enumerate}
    \item \textbf{Linearity}: $\curl(\vec{F} + \vec{G}) = \curl\vec{F} + \curl\vec{G}$
    \item \textbf{Product rule}: $\curl(\phi\vec{F}) = \phi\curl\vec{F} + \grad\phi \times \vec{F}$
    \item \textbf{Curl of gradient}: $\curl(\grad\phi) = \vec{0}$ (always zero!)
    \item \textbf{Divergence of curl}: $\div(\curl\vec{F}) = 0$ (always zero!)
\end{enumerate}

\subsubsection{Curl in Other Coordinate Systems}
\begin{formulabox}
\textbf{Cylindrical coordinates:}
\begin{equation}
\curl \vec{F} = \frac{1}{r}\begin{vmatrix}
\hat{r} & r\hat{\theta} & \hat{z} \\
\pdv{}{r} & \pdv{}{\theta} & \pdv{}{z} \\
F_r & rF_\theta & F_z
\end{vmatrix}
\end{equation}

\textbf{Spherical coordinates:}
\begin{equation}
\curl \vec{F} = \frac{1}{r^2\sin\theta}\begin{vmatrix}
\hat{r} & r\hat{\theta} & r\sin\theta\hat{\phi} \\
\pdv{}{r} & \pdv{}{\theta} & \pdv{}{\phi} \\
F_r & rF_\theta & r\sin\theta F_\phi
\end{vmatrix}
\end{equation}
\end{formulabox}

\begin{examplebox}
\textbf{Example:} Find the curl of $\vec{F} = y\hat{i} - x\hat{j}$.

\textbf{Solution:}
\begin{align}
\curl\vec{F} &= \begin{vmatrix}
\hat{i} & \hat{j} & \hat{k} \\
\pdv{}{x} & \pdv{}{y} & \pdv{}{z} \\
y & -x & 0
\end{vmatrix} \\
&= \left(\pdv{0}{y} - \pdv{(-x)}{z}\right)\hat{i} - \left(\pdv{0}{x} - \pdv{y}{z}\right)\hat{j} + \left(\pdv{(-x)}{x} - \pdv{y}{y}\right)\hat{k} \\
&= (0 - 0)\hat{i} - (0 - 0)\hat{j} + (-1 - 1)\hat{k} = -2\hat{k}
\end{align}
\end{examplebox}

\subsection{Laplacian}
The Laplacian is the divergence of the gradient of a scalar field.

\begin{formulabox}
\textbf{Laplacian in Cartesian coordinates:}
\begin{equation}
\laplacian \phi = \div(\grad\phi) = \pdv[2]{\phi}{x} + \pdv[2]{\phi}{y} + \pdv[2]{\phi}{z}
\end{equation}

\textbf{Vector Laplacian:}
\begin{equation}
\laplacian \vec{F} = \grad(\div\vec{F}) - \curl(\curl\vec{F})
\end{equation}
\end{formulabox}

\subsection{Line Integrals}
Line integrals integrate a function along a curve.

\subsubsection{Line Integral of Scalar Function}
\begin{formulabox}
\textbf{Line integral of scalar function:}
\begin{equation}
\int_C f(x,y,z) \, ds = \int_a^b f(x(t), y(t), z(t)) \sqrt{\left(\dv{x}{t}\right)^2 + \left(\dv{y}{t}\right)^2 + \left(\dv{z}{t}\right)^2} \, dt
\end{equation}

where $C$ is parameterized by $\vec{r}(t) = x(t)\hat{i} + y(t)\hat{j} + z(t)\hat{k}$, $a \leq t \leq b$.
\end{formulabox}

\subsubsection{Line Integral of Vector Field (Work)}
\begin{formulabox}
\textbf{Line integral of vector field (work done):}
\begin{equation}
\int_C \vec{F} \cdot d\vec{r} = \int_a^b \vec{F}(\vec{r}(t)) \cdot \dv{\vec{r}}{t} \, dt
\end{equation}

\textbf{Alternative forms:}
\begin{align}
\int_C \vec{F} \cdot d\vec{r} &= \int_C F_x dx + F_y dy + F_z dz \\
&= \int_C |\vec{F}| \cos\theta \, ds
\end{align}
\end{formulabox}

\subsubsection{Fundamental Theorem for Line Integrals}
\begin{importantbox}
If $\vec{F} = \grad\phi$ (conservative field), then:
\begin{equation}
\int_C \vec{F} \cdot d\vec{r} = \phi(\vec{r}(b)) - \phi(\vec{r}(a))
\end{equation}

The integral depends only on endpoints, not the path!
\end{importantbox}

\subsubsection{Path Independence}
A vector field $\vec{F}$ is conservative if and only if:
\begin{enumerate}
    \item $\curl\vec{F} = \vec{0}$ (in simply connected region)
    \item $\int_C \vec{F} \cdot d\vec{r} = 0$ for all closed curves $C$
    \item $\vec{F} = \grad\phi$ for some scalar potential $\phi$
\end{enumerate}

\begin{examplebox}
\textbf{Example:} Evaluate $\int_C (x^2 + y^2) \, ds$ where $C$ is the circle $x^2 + y^2 = 4$ from $(2,0)$ to $(0,2)$.

\textbf{Solution:} Parameterize: $x = 2\cos t$, $y = 2\sin t$, $0 \leq t \leq \pi/2$.

\begin{align}
ds &= \sqrt{(-2\sin t)^2 + (2\cos t)^2} \, dt = 2 \, dt \\
\int_C (x^2 + y^2) \, ds &= \int_0^{\pi/2} (4\cos^2 t + 4\sin^2 t) \cdot 2 \, dt \\
&= \int_0^{\pi/2} 4 \cdot 2 \, dt = 8 \cdot \frac{\pi}{2} = 4\pi
\end{align}
\end{examplebox}

\subsection{Surface Integrals}
Surface integrals integrate over a surface.

\subsubsection{Surface Integral of Scalar Function}
\begin{formulabox}
\textbf{Surface integral of scalar function:}
\begin{equation}
\iint_S f(x,y,z) \, dS = \iint_D f(\vec{r}(u,v)) \left|\pdv{\vec{r}}{u} \times \pdv{\vec{r}}{v}\right| \, du \, dv
\end{equation}

where $S$ is parameterized by $\vec{r}(u,v)$, and $D$ is the domain in $uv$-plane.
\end{formulabox}

\subsubsection{Surface Integral of Vector Field (Flux)}
\begin{formulabox}
\textbf{Surface integral of vector field (flux):}
\begin{equation}
\iint_S \vec{F} \cdot d\vec{S} = \iint_S \vec{F} \cdot \hat{n} \, dS
\end{equation}

where $\hat{n}$ is the unit normal vector to the surface.

\textbf{For $z = g(x,y)$:}
\begin{equation}
\iint_S \vec{F} \cdot d\vec{S} = \iint_D \left(-F_x \pdv{g}{x} - F_y \pdv{g}{y} + F_z\right) \, dx \, dy
\end{equation}
\end{formulabox}

\subsubsection{Orientation}
\begin{itemize}
    \item \textbf{Outward normal}: Points away from enclosed volume
    \item \textbf{Inward normal}: Points into enclosed volume
    \item For closed surfaces, convention is outward normal
\end{itemize}

\begin{examplebox}
\textbf{Example:} Find the flux of $\vec{F} = z\hat{k}$ through the surface $z = \sqrt{x^2 + y^2}$, $0 \leq z \leq 1$.

\textbf{Solution:} In cylindrical coordinates, $z = r$, $0 \leq r \leq 1$, $0 \leq \theta \leq 2\pi$.

The upward normal: $\hat{n} = \frac{-\pdv{g}{x}\hat{i} - \pdv{g}{y}\hat{j} + \hat{k}}{\sqrt{1 + (\pdv{g}{x})^2 + (\pdv{g}{y})^2}}$

For $z = \sqrt{x^2 + y^2} = r$:
\begin{align}
\pdv{g}{x} &= \frac{x}{\sqrt{x^2 + y^2}} = \cos\theta \\
\pdv{g}{y} &= \frac{y}{\sqrt{x^2 + y^2}} = \sin\theta
\end{align}

\begin{align}
\iint_S \vec{F} \cdot d\vec{S} &= \iint_D z \, dx \, dy = \iint_D r \cdot r \, dr \, d\theta \\
&= \int_0^{2\pi} \int_0^1 r^2 \, dr \, d\theta = 2\pi \cdot \frac{1}{3} = \frac{2\pi}{3}
\end{align}
\end{examplebox}

\subsection{Volume Integrals}
\begin{formulabox}
\textbf{Volume integral:}
\begin{equation}
\iiint_V f(x,y,z) \, dV = \iiint_V f(x,y,z) \, dx \, dy \, dz
\end{equation}
\end{formulabox}

\subsection{Important Theorems}
\subsubsection{Divergence Theorem (Gauss's Theorem)}
\begin{importantbox}
\textbf{Divergence Theorem:}
\begin{equation}
\iiint_V (\div\vec{F}) \, dV = \oiint_S \vec{F} \cdot d\vec{S}
\end{equation}

The volume integral of divergence equals the flux through the closed surface.
\end{importantbox}

\subsubsection{Stokes' Theorem}
\begin{importantbox}
\textbf{Stokes' Theorem:}
\begin{equation}
\iint_S (\curl\vec{F}) \cdot d\vec{S} = \oint_C \vec{F} \cdot d\vec{r}
\end{equation}

The surface integral of curl equals the line integral around the boundary.
\end{importantbox}

\subsubsection{Green's Theorem (2D)}
\begin{importantbox}
\textbf{Green's Theorem:}
\begin{equation}
\iint_D \left(\pdv{Q}{x} - \pdv{P}{y}\right) \, dx \, dy = \oint_C P \, dx + Q \, dy
\end{equation}
\end{importantbox}

\subsection{Practice Problems}
\begin{enumerate}
    \item Find $\grad\phi$ for $\phi = e^{xyz}$
    \item Show that $\div(\curl\vec{F}) = 0$ for any vector field
    \item Evaluate $\int_C (x^2 + y) \, dx + (y^2 + x) \, dy$ where $C$ is the line from $(0,0)$ to $(1,1)$
    \item Use divergence theorem to find flux of $\vec{F} = x^3\hat{i} + y^3\hat{j} + z^3\hat{k}$ through sphere $x^2 + y^2 + z^2 = a^2$
\end{enumerate}

\newpage

% ============================================
% COMPLEX ANALYSIS
% ============================================
\section{Complex Analysis}
Complex analysis is essential for solving integrals, understanding analytic functions, and working with complex potentials.

\subsection{Complex Numbers}
\subsubsection{Basic Definitions}
\begin{formulabox}
\textbf{Complex number:} $z = x + iy$ where $i^2 = -1$

\textbf{Real and imaginary parts:}
\begin{align}
\text{Re}(z) &= x \\
\text{Im}(z) &= y
\end{align}

\textbf{Polar form:}
\begin{align}
z &= r(\cos\theta + i\sin\theta) = re^{i\theta} \\
r &= |z| = \sqrt{x^2 + y^2} \\
\theta &= \arg(z) = \arctan\left(\frac{y}{x}\right)
\end{align}
\end{formulabox}

\subsubsection{Properties}
\begin{itemize}
    \item \textbf{Conjugate}: $\bar{z} = x - iy$, $z\bar{z} = |z|^2$
    \item \textbf{Modulus}: $|z| = \sqrt{z\bar{z}} = \sqrt{x^2 + y^2}$
    \item \textbf{Argument}: $\arg(z_1 z_2) = \arg(z_1) + \arg(z_2)$
    \item \textbf{De Moivre's theorem}: $(\cos\theta + i\sin\theta)^n = \cos(n\theta) + i\sin(n\theta)$
    \item \textbf{Euler's formula}: $e^{i\theta} = \cos\theta + i\sin\theta$
\end{itemize}

\subsection{Complex Functions}
\subsubsection{Definition}
A complex function $f(z)$ maps complex numbers to complex numbers:
\begin{equation}
f(z) = u(x,y) + iv(x,y)
\end{equation}
where $u$ and $v$ are real-valued functions.

\subsubsection{Limits and Continuity}
\begin{formulabox}
\textbf{Limit:} $\lim_{z \to z_0} f(z) = L$ if for every $\epsilon > 0$, there exists $\delta > 0$ such that $|f(z) - L| < \epsilon$ whenever $0 < |z - z_0| < \delta$.

\textbf{Continuity:} $f$ is continuous at $z_0$ if $\lim_{z \to z_0} f(z) = f(z_0)$.
\end{formulabox}

\subsection{Analytic Functions}
\subsubsection{Cauchy-Riemann Equations}
\begin{importantbox}
\textbf{Cauchy-Riemann Equations:}

A function $f(z) = u(x,y) + iv(x,y)$ is analytic if and only if:
\begin{align}
\pdv{u}{x} &= \pdv{v}{y} \\
\pdv{u}{y} &= -\pdv{v}{x}
\end{align}

These are necessary and sufficient conditions (in a domain) for differentiability.
\end{importantbox}

\subsubsection{Derivative of Complex Function}
\begin{formulabox}
\textbf{Complex derivative:}
\begin{equation}
f'(z) = \lim_{\Delta z \to 0} \frac{f(z + \Delta z) - f(z)}{\Delta z}
\end{equation}

If $f$ is analytic, then:
\begin{equation}
f'(z) = \pdv{u}{x} + i\pdv{v}{x} = \pdv{v}{y} - i\pdv{u}{y}
\end{equation}
\end{formulabox}

\subsubsection{Harmonic Functions}
If $f = u + iv$ is analytic, then $u$ and $v$ are harmonic:
\begin{equation}
\pdv[2]{u}{x} + \pdv[2]{u}{y} = 0, \quad \pdv[2]{v}{x} + \pdv[2]{v}{y} = 0
\end{equation}

\begin{examplebox}
\textbf{Example:} Show that $f(z) = z^2$ is analytic.

\textbf{Solution:} $f(z) = (x + iy)^2 = x^2 - y^2 + i(2xy)$

So $u = x^2 - y^2$ and $v = 2xy$.

Check Cauchy-Riemann:
\begin{align}
\pdv{u}{x} &= 2x = \pdv{v}{y} = 2x \quad \checkmark \\
\pdv{u}{y} &= -2y = -\pdv{v}{x} = -2y \quad \checkmark
\end{align}

Therefore, $f(z)$ is analytic and $f'(z) = 2z$.
\end{examplebox}

\subsection{Complex Integration}
\subsubsection{Line Integral in Complex Plane}
\begin{formulabox}
\textbf{Complex line integral:}
\begin{equation}
\int_C f(z) \, dz = \int_a^b f(z(t)) \dv{z}{t} \, dt
\end{equation}

where $C$ is parameterized by $z(t)$, $a \leq t \leq b$.
\end{formulabox}

\subsubsection{Properties}
\begin{enumerate}
    \item \textbf{Linearity}: $\int_C [\alpha f(z) + \beta g(z)] \, dz = \alpha\int_C f(z) \, dz + \beta\int_C g(z) \, dz$
    \item \textbf{Path reversal}: $\int_{-C} f(z) \, dz = -\int_C f(z) \, dz$
    \item \textbf{Path addition}: $\int_{C_1 + C_2} f(z) \, dz = \int_{C_1} f(z) \, dz + \int_{C_2} f(z) \, dz$
    \item \textbf{ML-inequality}: $\left|\int_C f(z) \, dz\right| \leq ML$ where $M = \max|f(z)|$ on $C$ and $L$ is the length of $C$
\end{enumerate}

\subsection{Cauchy's Theorem}
\begin{importantbox}
\textbf{Cauchy's Theorem (Fundamental):}

If $f(z)$ is analytic in a simply connected domain $D$ and $C$ is a closed contour in $D$, then:
\begin{equation}
\oint_C f(z) \, dz = 0
\end{equation}

\textbf{Consequence:} The integral is path-independent in simply connected domains!
\end{importantbox}

\subsubsection{Cauchy's Integral Formula}
\begin{importantbox}
\textbf{Cauchy's Integral Formula:}

If $f(z)$ is analytic inside and on a simple closed contour $C$, and $z_0$ is inside $C$, then:
\begin{equation}
f(z_0) = \frac{1}{2\pi i} \oint_C \frac{f(z)}{z - z_0} \, dz
\end{equation}

\textbf{Generalized form:}
\begin{equation}
f^{(n)}(z_0) = \frac{n!}{2\pi i} \oint_C \frac{f(z)}{(z - z_0)^{n+1}} \, dz
\end{equation}
\end{importantbox}

\begin{examplebox}
\textbf{Example:} Evaluate $\oint_C \frac{e^z}{z^2} \, dz$ where $C$ is the circle $|z| = 1$.

\textbf{Solution:} Using Cauchy's integral formula with $f(z) = e^z$ and $z_0 = 0$:

\begin{align}
f'(0) &= \frac{1!}{2\pi i} \oint_C \frac{e^z}{(z - 0)^2} \, dz \\
e^0 = 1 &= \frac{1}{2\pi i} \oint_C \frac{e^z}{z^2} \, dz \\
\oint_C \frac{e^z}{z^2} \, dz &= 2\pi i
\end{align}
\end{examplebox}

\subsection{Laurent Series}
\subsubsection{Definition}
\begin{formulabox}
\textbf{Laurent series expansion:}

For a function analytic in an annulus $R_1 < |z - z_0| < R_2$:
\begin{equation}
f(z) = \sum_{n=-\infty}^{\infty} a_n (z - z_0)^n
\end{equation}

where
\begin{equation}
a_n = \frac{1}{2\pi i} \oint_C \frac{f(z)}{(z - z_0)^{n+1}} \, dz
\end{equation}
\end{formulabox}

\subsubsection{Classification of Singularities}
\begin{enumerate}
    \item \textbf{Removable singularity}: All $a_n = 0$ for $n < 0$
    \item \textbf{Pole of order $m$}: $a_{-m} \neq 0$ but $a_n = 0$ for $n < -m$
    \item \textbf{Essential singularity}: Infinitely many nonzero $a_n$ for $n < 0$
\end{enumerate}

\subsection{Residue Calculus}
\subsubsection{Residue}
\begin{formulabox}
\textbf{Residue at isolated singularity:}

The residue of $f(z)$ at $z_0$ is the coefficient $a_{-1}$ in the Laurent expansion:
\begin{equation}
\text{Res}(f, z_0) = a_{-1} = \frac{1}{2\pi i} \oint_C f(z) \, dz
\end{equation}
\end{formulabox}

\subsubsection{Computing Residues}
\begin{formulabox}
\textbf{Simple pole:} If $f(z) = \frac{g(z)}{h(z)}$ where $h(z_0) = 0$, $h'(z_0) \neq 0$, and $g(z_0) \neq 0$:
\begin{equation}
\text{Res}(f, z_0) = \frac{g(z_0)}{h'(z_0)}
\end{equation}

\textbf{Pole of order $m$:} If $f(z)$ has a pole of order $m$ at $z_0$:
\begin{equation}
\text{Res}(f, z_0) = \frac{1}{(m-1)!} \lim_{z \to z_0} \dv[m-1]{}{z}\left[(z - z_0)^m f(z)\right]
\end{equation}
\end{formulabox}

\subsubsection{Residue Theorem}
\begin{importantbox}
\textbf{Residue Theorem:}

If $f(z)$ is analytic inside and on a simple closed contour $C$ except for isolated singularities $z_1, z_2, \ldots, z_n$ inside $C$, then:
\begin{equation}
\oint_C f(z) \, dz = 2\pi i \sum_{k=1}^n \text{Res}(f, z_k)
\end{equation}
\end{importantbox}

\begin{examplebox}
\textbf{Example:} Evaluate $\oint_C \frac{1}{z^2 + 1} \, dz$ where $C$ is the circle $|z| = 2$.

\textbf{Solution:} The singularities are at $z = \pm i$, both inside $C$.

For $z = i$:
\begin{align}
\text{Res}\left(\frac{1}{z^2 + 1}, i\right) &= \frac{1}{2i} = -\frac{i}{2}
\end{align}

For $z = -i$:
\begin{align}
\text{Res}\left(\frac{1}{z^2 + 1}, -i\right) &= \frac{1}{-2i} = \frac{i}{2}
\end{align}

\begin{align}
\oint_C \frac{1}{z^2 + 1} \, dz &= 2\pi i\left(-\frac{i}{2} + \frac{i}{2}\right) = 0
\end{align}
\end{examplebox}

\subsection{Contour Integration}
Contour integration is a powerful technique for evaluating real integrals.

\subsubsection{Types of Contours}
\begin{enumerate}
    \item \textbf{Semicircular contour}: For integrals $\int_{-\infty}^{\infty} f(x) \, dx$
    \item \textbf{Keyhole contour}: For integrals with branch cuts
    \item \textbf{Rectangular contour}: For periodic functions
\end{enumerate}

\subsubsection{Standard Integrals}
\begin{formulabox}
\textbf{Integral of rational functions:}
\begin{equation}
\int_{-\infty}^{\infty} \frac{P(x)}{Q(x)} \, dx = 2\pi i \sum_{\text{upper half-plane}} \text{Res}\left(\frac{P(z)}{Q(z)}, z_k\right)
\end{equation}

provided $\deg Q \geq \deg P + 2$ and no poles on real axis.

\textbf{Integrals with $e^{iax}$:}
\begin{equation}
\int_{-\infty}^{\infty} \frac{e^{iax}}{Q(x)} \, dx = 2\pi i \sum_{\text{upper half-plane}} \text{Res}\left(\frac{e^{iaz}}{Q(z)}, z_k\right)
\end{equation}
\end{formulabox}

\begin{examplebox}
\textbf{Example:} Evaluate $I = \int_{-\infty}^{\infty} \frac{1}{x^2 + 1} \, dx$.

\textbf{Solution:} Consider $f(z) = \frac{1}{z^2 + 1}$. The only pole in the upper half-plane is $z = i$.

\begin{align}
\text{Res}(f, i) &= \lim_{z \to i} (z - i) \frac{1}{(z-i)(z+i)} = \frac{1}{2i}
\end{align}

Using residue theorem on semicircular contour:
\begin{align}
I &= 2\pi i \cdot \frac{1}{2i} = \pi
\end{align}
\end{examplebox}

\subsubsection{Jordan's Lemma}
\begin{importantbox}
\textbf{Jordan's Lemma:}

If $f(z) \to 0$ uniformly as $|z| \to \infty$ in the upper half-plane, then:
\begin{equation}
\lim_{R \to \infty} \int_{C_R} e^{iaz} f(z) \, dz = 0 \quad (a > 0)
\end{equation}

where $C_R$ is the semicircle $|z| = R$, $\text{Im}(z) \geq 0$.
\end{importantbox}

\subsection{Practice Problems}
\begin{enumerate}
    \item Show that $f(z) = \bar{z}$ is not analytic
    \item Find all singularities of $f(z) = \frac{1}{\sin z}$
    \item Evaluate $\oint_C \frac{\cos z}{z^3} \, dz$ where $C$ is $|z| = 1$
    \item Use residue theorem to evaluate $\int_{-\infty}^{\infty} \frac{x^2}{x^4 + 1} \, dx$
\end{enumerate}

\newpage

% ============================================
% DIFFERENTIAL EQUATIONS
% ============================================
\section{Differential Equations}
Differential equations are fundamental in physics for describing dynamics and field equations.

\subsection{First-Order ODEs}
\subsubsection{Separable Equations}
\begin{formulabox}
\textbf{Separable ODE:}
\begin{equation}
\dv{y}{x} = f(x)g(y)
\end{equation}

Solution: Separate and integrate:
\begin{equation}
\int \frac{dy}{g(y)} = \int f(x) \, dx
\end{equation}
\end{formulabox}

\subsubsection{Linear First-Order ODE}
\begin{formulabox}
\textbf{Linear ODE:}
\begin{equation}
\dv{y}{x} + P(x)y = Q(x)
\end{equation}

\textbf{Integrating factor:}
\begin{equation}
\mu(x) = e^{\int P(x) \, dx}
\end{equation}

\textbf{Solution:}
\begin{equation}
y(x) = \frac{1}{\mu(x)}\left[\int \mu(x)Q(x) \, dx + C\right]
\end{equation}
\end{formulabox}

\subsubsection{Exact Equations}
\begin{formulabox}
\textbf{Exact ODE:}
\begin{equation}
M(x,y) \, dx + N(x,y) \, dy = 0
\end{equation}

is exact if $\pdv{M}{y} = \pdv{N}{x}$.

Solution: Find $\phi(x,y)$ such that $d\phi = M \, dx + N \, dy = 0$, so $\phi(x,y) = C$.
\end{formulabox}

\subsection{Second-Order Linear ODEs}
\subsubsection{General Form}
\begin{formulabox}
\textbf{Second-order linear ODE:}
\begin{equation}
a(x)\dv[2]{y}{x} + b(x)\dv{y}{x} + c(x)y = f(x)
\end{equation}

If $f(x) = 0$, it's homogeneous; otherwise, it's non-homogeneous.
\end{formulabox}

\subsubsection{Constant Coefficients}
\begin{formulabox}
\textbf{Homogeneous with constant coefficients:}
\begin{equation}
a\dv[2]{y}{x} + b\dv{y}{x} + cy = 0
\end{equation}

\textbf{Characteristic equation:}
\begin{equation}
ar^2 + br + c = 0
\end{equation}

\textbf{Solutions:}
\begin{itemize}
    \item \textbf{Distinct real roots} $r_1, r_2$: $y = C_1 e^{r_1 x} + C_2 e^{r_2 x}$
    \item \textbf{Repeated root} $r$: $y = (C_1 + C_2 x)e^{rx}$
    \item \textbf{Complex roots} $\alpha \pm i\beta$: $y = e^{\alpha x}(C_1 \cos\beta x + C_2 \sin\beta x)$
\end{itemize}
\end{formulabox}

\subsection{Series Solutions: Frobenius Method}
\subsubsection{Ordinary Points}
\begin{formulabox}
\textbf{Power series solution about ordinary point:}

For ODE: $P(x)\dv[2]{y}{x} + Q(x)\dv{y}{x} + R(x)y = 0$

If $x_0$ is an ordinary point (not a singularity), solution is:
\begin{equation}
y(x) = \sum_{n=0}^{\infty} a_n (x - x_0)^n
\end{equation}
\end{formulabox}

\subsubsection{Regular Singular Points}
\begin{formulabox}
\textbf{Regular singular point:}

$x_0$ is a regular singular point if $(x-x_0)P(x)$ and $(x-x_0)^2 Q(x)$ are analytic at $x_0$.

\textbf{Indicial equation:} For ODE in form:
\begin{equation}
x^2\dv[2]{y}{x} + xp(x)\dv{y}{x} + q(x)y = 0
\end{equation}

The indicial equation is:
\begin{equation}
r(r-1) + p_0 r + q_0 = 0
\end{equation}

where $p_0 = \lim_{x \to 0} xp(x)$ and $q_0 = \lim_{x \to 0} x^2 q(x)$.
\end{formulabox}

\subsubsection{Frobenius Method}
\begin{importantbox}
\textbf{Frobenius Method:}

For ODE with regular singular point at $x = 0$:
\begin{equation}
x^2\dv[2]{y}{x} + xp(x)\dv{y}{x} + q(x)y = 0
\end{equation}

Assume solution:
\begin{equation}
y(x) = x^r \sum_{n=0}^{\infty} a_n x^n = \sum_{n=0}^{\infty} a_n x^{n+r}
\end{equation}

Substitute into ODE, find indicial equation for $r$, then determine recurrence relation for $a_n$.
\end{importantbox}

\subsubsection{Cases of Frobenius Method}
\begin{enumerate}
    \item \textbf{Case 1:} $r_1 - r_2 \notin \mathbb{Z}$: Two independent solutions
    \item \textbf{Case 2:} $r_1 = r_2$: Second solution involves logarithm
    \item \textbf{Case 3:} $r_1 - r_2 \in \mathbb{Z}$: May need logarithmic term
\end{enumerate}

\begin{examplebox}
\textbf{Example:} Solve Bessel's equation of order 0: $x^2\dv[2]{y}{x} + x\dv{y}{x} + x^2 y = 0$.

\textbf{Solution:} $x = 0$ is a regular singular point.

Indicial equation: $r^2 = 0$, so $r = 0$ (double root).

Using Frobenius method with $y = \sum_{n=0}^{\infty} a_n x^n$:

Recurrence: $a_n = -\frac{a_{n-2}}{n^2}$ for $n \geq 2$, with $a_1 = 0$.

Solution: $J_0(x) = \sum_{m=0}^{\infty} \frac{(-1)^m}{(m!)^2}\left(\frac{x}{2}\right)^{2m}$ (Bessel function of first kind).
\end{examplebox}

\subsection{Special Functions}
\subsubsection{Bessel Functions}
\begin{formulabox}
\textbf{Bessel's equation of order $\nu$:}
\begin{equation}
x^2\dv[2]{y}{x} + x\dv{y}{x} + (x^2 - \nu^2)y = 0
\end{equation}

\textbf{Bessel function of first kind:}
\begin{equation}
J_\nu(x) = \sum_{m=0}^{\infty} \frac{(-1)^m}{m!\Gamma(m+\nu+1)}\left(\frac{x}{2}\right)^{2m+\nu}
\end{equation}

\textbf{Bessel function of second kind:}
\begin{equation}
Y_\nu(x) = \frac{J_\nu(x)\cos(\nu\pi) - J_{-\nu}(x)}{\sin(\nu\pi)}
\end{equation}
\end{formulabox}

\subsubsection{Properties of Bessel Functions}
\begin{itemize}
    \item \textbf{Recurrence relations:}
    \begin{align}
    J_{\nu-1}(x) + J_{\nu+1}(x) &= \frac{2\nu}{x}J_\nu(x) \\
    J_{\nu-1}(x) - J_{\nu+1}(x) &= 2J_\nu'(x)
    \end{align}
    \item \textbf{Orthogonality:}
    \begin{equation}
    \int_0^1 x J_\nu(\alpha_m x) J_\nu(\alpha_n x) \, dx = 0 \quad (m \neq n)
    \end{equation}
    where $\alpha_m$ are zeros of $J_\nu$.
\end{itemize}

\subsubsection{Legendre Polynomials}
\begin{formulabox}
\textbf{Legendre's equation:}
\begin{equation}
(1-x^2)\dv[2]{y}{x} - 2x\dv{y}{x} + n(n+1)y = 0
\end{equation}

\textbf{Legendre polynomials (Rodrigues formula):}
\begin{equation}
P_n(x) = \frac{1}{2^n n!}\dv[n]{}{x}[(x^2-1)^n]
\end{equation}

\textbf{First few:}
\begin{align}
P_0(x) &= 1 \\
P_1(x) &= x \\
P_2(x) &= \frac{1}{2}(3x^2 - 1) \\
P_3(x) &= \frac{1}{2}(5x^3 - 3x)
\end{align}
\end{formulabox}

\subsubsection{Properties of Legendre Polynomials}
\begin{itemize}
    \item \textbf{Orthogonality:}
    \begin{equation}
    \int_{-1}^1 P_m(x) P_n(x) \, dx = \frac{2}{2n+1}\delta_{mn}
    \end{equation}
    \item \textbf{Generating function:}
    \begin{equation}
    \frac{1}{\sqrt{1-2xt+t^2}} = \sum_{n=0}^{\infty} P_n(x) t^n
    \end{equation}
    \item \textbf{Recurrence:}
    \begin{equation}
    (n+1)P_{n+1}(x) = (2n+1)xP_n(x) - nP_{n-1}(x)
    \end{equation}
\end{itemize}

\subsubsection{Hermite Polynomials}
\begin{formulabox}
\textbf{Hermite's equation:}
\begin{equation}
\dv[2]{y}{x} - 2x\dv{y}{x} + 2ny = 0
\end{equation}

\textbf{Hermite polynomials (Rodrigues formula):}
\begin{equation}
H_n(x) = (-1)^n e^{x^2} \dv[n]{}{x}(e^{-x^2})
\end{equation}

\textbf{First few:}
\begin{align}
H_0(x) &= 1 \\
H_1(x) &= 2x \\
H_2(x) &= 4x^2 - 2 \\
H_3(x) &= 8x^3 - 12x
\end{align}
\end{formulabox}

\subsubsection{Properties of Hermite Polynomials}
\begin{itemize}
    \item \textbf{Orthogonality:}
    \begin{equation}
    \int_{-\infty}^{\infty} H_m(x) H_n(x) e^{-x^2} \, dx = 2^n n!\sqrt{\pi}\delta_{mn}
    \end{equation}
    \item \textbf{Recurrence:}
    \begin{equation}
    H_{n+1}(x) = 2xH_n(x) - 2nH_{n-1}(x)
    \end{equation}
\end{itemize}

\subsubsection{Laguerre Polynomials}
\begin{formulabox}
\textbf{Laguerre's equation:}
\begin{equation}
x\dv[2]{y}{x} + (1-x)\dv{y}{x} + ny = 0
\end{equation}

\textbf{Laguerre polynomials:}
\begin{equation}
L_n(x) = \frac{e^x}{n!}\dv[n]{}{x}(x^n e^{-x})
\end{equation}
\end{formulabox}

\subsubsection{Gamma Function}
\begin{formulabox}
\textbf{Gamma function:}
\begin{equation}
\Gamma(z) = \int_0^{\infty} t^{z-1} e^{-t} \, dt \quad (\text{Re}(z) > 0)
\end{equation}

\textbf{Properties:}
\begin{align}
\Gamma(n+1) &= n! \quad (n \in \mathbb{N}) \\
\Gamma(z+1) &= z\Gamma(z) \\
\Gamma(1/2) &= \sqrt{\pi} \\
\Gamma(z)\Gamma(1-z) &= \frac{\pi}{\sin(\pi z)}
\end{align}
\end{formulabox}

\subsection{Sturm-Liouville Theory}
\subsubsection{Sturm-Liouville Problem}
\begin{formulabox}
\textbf{Sturm-Liouville equation:}
\begin{equation}
\dv{}{x}\left[p(x)\dv{y}{x}\right] + [q(x) + \lambda w(x)]y = 0
\end{equation}

with boundary conditions. The eigenvalues $\lambda_n$ are real and eigenfunctions $y_n(x)$ are orthogonal with weight $w(x)$.
\end{formulabox}

\subsection{Practice Problems}
\begin{enumerate}
    \item Solve $x^2\dv[2]{y}{x} + x\dv{y}{x} + (x^2 - 4)y = 0$ using Frobenius method
    \item Find the first three Legendre polynomials using Rodrigues formula
    \item Show that $J_0'(x) = -J_1(x)$
    \item Solve the Hermite equation for $n = 2$
\end{enumerate}

\newpage

% ============================================
% LINEAR ALGEBRA
% ============================================
\section{Linear Algebra}
Linear algebra is essential for quantum mechanics, classical mechanics, and many-body systems.

\subsection{Vector Spaces}
\subsubsection{Definition}
\begin{formulabox}
\textbf{Vector space} $V$ over field $\mathbb{F}$ satisfies:
\begin{enumerate}
    \item Closure under addition: $\vec{u} + \vec{v} \in V$
    \item Closure under scalar multiplication: $\alpha\vec{u} \in V$
    \item Commutativity: $\vec{u} + \vec{v} = \vec{v} + \vec{u}$
    \item Associativity: $(\vec{u} + \vec{v}) + \vec{w} = \vec{u} + (\vec{v} + \vec{w})$
    \item Identity: $\vec{0} + \vec{u} = \vec{u}$
    \item Inverse: $\vec{u} + (-\vec{u}) = \vec{0}$
    \item Distributivity: $\alpha(\vec{u} + \vec{v}) = \alpha\vec{u} + \alpha\vec{v}$
\end{enumerate}
\end{formulabox}

\subsubsection{Linear Independence}
\begin{formulabox}
\textbf{Linear independence:}

Vectors $\{\vec{v}_1, \vec{v}_2, \ldots, \vec{v}_n\}$ are linearly independent if:
\begin{equation}
\sum_{i=1}^n c_i \vec{v}_i = \vec{0} \quad \Rightarrow \quad c_i = 0 \text{ for all } i
\end{equation}

Otherwise, they are linearly dependent.
\end{formulabox}

\subsubsection{Basis and Dimension}
\begin{formulabox}
\textbf{Basis:} A set of linearly independent vectors that span the vector space.

\textbf{Dimension:} The number of vectors in a basis (unique for finite-dimensional spaces).
\end{formulabox}

\subsection{Matrices}
\subsubsection{Basic Operations}
\begin{formulabox}
\textbf{Matrix multiplication:} $(AB)_{ij} = \sum_k A_{ik} B_{kj}$

\textbf{Transpose:} $(A^T)_{ij} = A_{ji}$

\textbf{Conjugate transpose:} $(A^\dagger)_{ij} = \bar{A}_{ji}$

\textbf{Inverse:} $AA^{-1} = A^{-1}A = I$
\end{formulabox}

\subsubsection{Determinant}
\begin{formulabox}
\textbf{Determinant (2×2):}
\begin{equation}
\det\begin{pmatrix} a & b \\ c & d \end{pmatrix} = ad - bc
\end{equation}

\textbf{Determinant (3×3):}
\begin{equation}
\det\begin{pmatrix} a & b & c \\ d & e & f \\ g & h & i \end{pmatrix} = a(ei - fh) - b(di - fg) + c(dh - eg)
\end{equation}

\textbf{Properties:}
\begin{itemize}
    \item $\det(AB) = \det(A)\det(B)$
    \item $\det(A^T) = \det(A)$
    \item $\det(A^{-1}) = 1/\det(A)$
    \item $\det(\alpha A) = \alpha^n \det(A)$ for $n \times n$ matrix
\end{itemize}
\end{formulabox}

\subsection{Eigenvalues and Eigenvectors}
\subsubsection{Definition}
\begin{importantbox}
\textbf{Eigenvalue problem:}

For matrix $A$, find $\lambda$ (eigenvalue) and $\vec{v} \neq \vec{0}$ (eigenvector) such that:
\begin{equation}
A\vec{v} = \lambda\vec{v}
\end{equation}

\textbf{Characteristic equation:}
\begin{equation}
\det(A - \lambda I) = 0
\end{equation}
\end{importantbox}

\subsubsection{Properties}
\begin{enumerate}
    \item Sum of eigenvalues = trace of $A$
    \item Product of eigenvalues = determinant of $A$
    \item Eigenvectors corresponding to distinct eigenvalues are linearly independent
    \item If $A$ is real and symmetric, eigenvalues are real
    \item If $A$ is Hermitian, eigenvalues are real
\end{enumerate}

\subsubsection{Multiplicity}
\begin{itemize}
    \item \textbf{Algebraic multiplicity}: Multiplicity of $\lambda$ as a root of characteristic polynomial
    \item \textbf{Geometric multiplicity}: Dimension of eigenspace $E_\lambda = \{\vec{v} : A\vec{v} = \lambda\vec{v}\}$
    \item Geometric multiplicity $\leq$ Algebraic multiplicity
\end{itemize}

\begin{examplebox}
\textbf{Example:} Find eigenvalues and eigenvectors of $A = \begin{pmatrix} 2 & 1 \\ 1 & 2 \end{pmatrix}$.

\textbf{Solution:} Characteristic equation:
\begin{align}
\det(A - \lambda I) &= \det\begin{pmatrix} 2-\lambda & 1 \\ 1 & 2-\lambda \end{pmatrix} \\
&= (2-\lambda)^2 - 1 = \lambda^2 - 4\lambda + 3 = 0
\end{align}

Eigenvalues: $\lambda_1 = 3$, $\lambda_2 = 1$.

For $\lambda_1 = 3$:
\begin{equation}
\begin{pmatrix} -1 & 1 \\ 1 & -1 \end{pmatrix}\begin{pmatrix} v_1 \\ v_2 \end{pmatrix} = \vec{0} \Rightarrow \vec{v}_1 = \begin{pmatrix} 1 \\ 1 \end{pmatrix}
\end{equation}

For $\lambda_2 = 1$:
\begin{equation}
\begin{pmatrix} 1 & 1 \\ 1 & 1 \end{pmatrix}\begin{pmatrix} v_1 \\ v_2 \end{pmatrix} = \vec{0} \Rightarrow \vec{v}_2 = \begin{pmatrix} 1 \\ -1 \end{pmatrix}
\end{equation}
\end{examplebox}

\subsection{Diagonalization}
\subsubsection{Definition}
\begin{formulabox}
\textbf{Diagonalization:}

A matrix $A$ is diagonalizable if there exists invertible matrix $P$ and diagonal matrix $D$ such that:
\begin{equation}
A = PDP^{-1} \quad \text{or} \quad D = P^{-1}AP
\end{equation}

where columns of $P$ are eigenvectors and $D$ has eigenvalues on diagonal.
\end{formulabox}

\subsubsection{Conditions for Diagonalization}
\begin{importantbox}
\textbf{Diagonalizable if and only if:}
\begin{enumerate}
    \item $A$ has $n$ linearly independent eigenvectors (for $n \times n$ matrix)
    \item Geometric multiplicity = Algebraic multiplicity for each eigenvalue
    \item Sum of geometric multiplicities = $n$
\end{enumerate}
\end{importantbox}

\subsubsection{Diagonalization Process}
\begin{enumerate}
    \item Find eigenvalues: Solve $\det(A - \lambda I) = 0$
    \item Find eigenvectors: Solve $(A - \lambda I)\vec{v} = \vec{0}$ for each $\lambda$
    \item Form $P$: Columns are eigenvectors
    \item Form $D$: Diagonal entries are eigenvalues (in same order)
    \item Verify: $A = PDP^{-1}$
\end{enumerate}

\begin{examplebox}
\textbf{Example:} Diagonalize $A = \begin{pmatrix} 2 & 1 \\ 1 & 2 \end{pmatrix}$.

\textbf{Solution:} From previous example, eigenvalues are $\lambda_1 = 3$, $\lambda_2 = 1$ with eigenvectors $\vec{v}_1 = \begin{pmatrix} 1 \\ 1 \end{pmatrix}$, $\vec{v}_2 = \begin{pmatrix} 1 \\ -1 \end{pmatrix}$.

\begin{align}
P &= \begin{pmatrix} 1 & 1 \\ 1 & -1 \end{pmatrix}, \quad P^{-1} = \frac{1}{2}\begin{pmatrix} 1 & 1 \\ 1 & -1 \end{pmatrix} \\
D &= \begin{pmatrix} 3 & 0 \\ 0 & 1 \end{pmatrix}
\end{align}

Check: $PDP^{-1} = A$ ✓
\end{examplebox}

\subsubsection{Applications}
\begin{itemize}
    \item \textbf{Matrix powers:} $A^n = PD^nP^{-1}$
    \item \textbf{Matrix exponentials:} $e^A = Pe^DP^{-1}$ where $e^D$ is diagonal with $e^{\lambda_i}$
    \item \textbf{Solving systems:} $\dv{\vec{x}}{t} = A\vec{x}$ has solution $\vec{x}(t) = e^{At}\vec{x}(0)$
\end{itemize}

\subsection{Hermitian Matrices}
\subsubsection{Definition}
\begin{formulabox}
\textbf{Hermitian (self-adjoint) matrix:}

$A$ is Hermitian if $A^\dagger = A$, i.e., $A_{ij} = \bar{A}_{ji}$.

For real matrices, this means $A$ is symmetric: $A^T = A$.
\end{formulabox}

\subsubsection{Properties of Hermitian Matrices}
\begin{importantbox}
\textbf{Key properties:}
\begin{enumerate}
    \item \textbf{Real eigenvalues}: All eigenvalues are real
    \item \textbf{Orthogonal eigenvectors}: Eigenvectors corresponding to distinct eigenvalues are orthogonal
    \item \textbf{Diagonalizable}: Always diagonalizable by unitary matrix
    \item \textbf{Spectral theorem}: $A = UDU^\dagger$ where $U$ is unitary and $D$ is real diagonal
    \item \textbf{Positive definite}: If all eigenvalues $> 0$
    \item \textbf{Positive semidefinite}: If all eigenvalues $\geq 0$
\end{enumerate}
\end{importantbox}

\subsubsection{Unitary Diagonalization}
\begin{formulabox}
\textbf{Unitary diagonalization:}

For Hermitian $A$, there exists unitary $U$ (satisfying $U^\dagger U = I$) such that:
\begin{equation}
A = UDU^\dagger
\end{equation}

where $D$ is real diagonal. Columns of $U$ are orthonormal eigenvectors.
\end{formulabox}

\subsubsection{Unitary Matrices}
\begin{formulabox}
\textbf{Unitary matrix:}

$U$ is unitary if $U^\dagger U = UU^\dagger = I$.

Properties:
\begin{itemize}
    \item Columns (and rows) are orthonormal
    \item $|\det(U)| = 1$
    \item Eigenvalues have unit modulus: $|\lambda| = 1$
    \item Preserves inner product: $(U\vec{x}) \cdot (U\vec{y}) = \vec{x} \cdot \vec{y}$
\end{itemize}
\end{formulabox}

\begin{examplebox}
\textbf{Example:} Show that $A = \begin{pmatrix} 1 & i \\ -i & 1 \end{pmatrix}$ is Hermitian and find its eigenvalues.

\textbf{Solution:} Check: $A^\dagger = \begin{pmatrix} 1 & i \\ -i & 1 \end{pmatrix}^T = \begin{pmatrix} 1 & -i \\ i & 1 \end{pmatrix}^* = \begin{pmatrix} 1 & i \\ -i & 1 \end{pmatrix} = A$ ✓

Characteristic equation:
\begin{align}
\det(A - \lambda I) &= \det\begin{pmatrix} 1-\lambda & i \\ -i & 1-\lambda \end{pmatrix} \\
&= (1-\lambda)^2 - (-i)(i) = (1-\lambda)^2 - 1 = \lambda^2 - 2\lambda = 0
\end{align}

Eigenvalues: $\lambda_1 = 0$, $\lambda_2 = 2$ (both real, as expected for Hermitian matrix).
\end{examplebox}

\subsection{Inner Product Spaces}
\subsubsection{Definition}
\begin{formulabox}
\textbf{Inner product:}

For vectors $\vec{u}, \vec{v}$ in complex vector space:
\begin{equation}
\langle \vec{u}, \vec{v} \rangle = \sum_i \bar{u}_i v_i = \vec{u}^\dagger \vec{v}
\end{equation}

Properties:
\begin{itemize}
    \item $\langle \vec{u}, \vec{v} \rangle = \overline{\langle \vec{v}, \vec{u} \rangle}$
    \item $\langle \alpha\vec{u} + \beta\vec{w}, \vec{v} \rangle = \alpha\langle \vec{u}, \vec{v} \rangle + \beta\langle \vec{w}, \vec{v} \rangle$
    \item $\langle \vec{u}, \vec{u} \rangle \geq 0$ with equality iff $\vec{u} = \vec{0}$
\end{itemize}
\end{formulabox}

\subsubsection{Orthogonality}
\begin{formulabox}
\textbf{Orthogonal vectors:} $\langle \vec{u}, \vec{v} \rangle = 0$

\textbf{Orthonormal basis:} $\{\vec{e}_i\}$ such that $\langle \vec{e}_i, \vec{e}_j \rangle = \delta_{ij}$

\textbf{Gram-Schmidt process:} Orthogonalize any basis to get orthonormal basis.
\end{formulabox}

\subsection{Matrix Decompositions}
\subsubsection{Singular Value Decomposition (SVD)}
\begin{formulabox}
\textbf{SVD:}

For any $m \times n$ matrix $A$:
\begin{equation}
A = U\Sigma V^\dagger
\end{equation}

where $U$ is $m \times m$ unitary, $V$ is $n \times n$ unitary, and $\Sigma$ is $m \times n$ diagonal with singular values $\sigma_i \geq 0$.
\end{formulabox}

\subsubsection{QR Decomposition}
\begin{formulabox}
\textbf{QR decomposition:}

For $m \times n$ matrix $A$ with linearly independent columns:
\begin{equation}
A = QR
\end{equation}

where $Q$ is $m \times n$ with orthonormal columns and $R$ is $n \times n$ upper triangular.
\end{formulabox}

\subsection{Applications in Physics}
\subsubsection{Quantum Mechanics}
\begin{itemize}
    \item Operators are Hermitian (observables have real eigenvalues)
    \item Eigenstates are orthogonal
    \item Unitary evolution: $|\psi(t)\rangle = U(t)|\psi(0)\rangle$
\end{itemize}

\subsubsection{Classical Mechanics}
\begin{itemize}
    \item Normal modes: Diagonalize mass and spring constant matrices
    \item Principal axes: Diagonalize moment of inertia tensor
\end{itemize}

\subsection{Practice Problems}
\begin{enumerate}
    \item Find eigenvalues and eigenvectors of $\begin{pmatrix} 3 & 1 \\ 1 & 3 \end{pmatrix}$
    \item Diagonalize the matrix in problem 1
    \item Show that eigenvalues of a Hermitian matrix are real
    \item Find a unitary matrix that diagonalizes $\begin{pmatrix} 2 & 1+i \\ 1-i & 3 \end{pmatrix}$
    \item Prove that eigenvectors of a Hermitian matrix with distinct eigenvalues are orthogonal
\end{enumerate}

\newpage

% ============================================
% SUMMARY AND QUICK REFERENCE
% ============================================
\section{Quick Reference and Summary}

\subsection{Vector Calculus Formulas}
\begin{multicols}{2}
\begin{itemize}
    \item $\grad\phi = \pdv{\phi}{x}\hat{i} + \pdv{\phi}{y}\hat{j} + \pdv{\phi}{z}\hat{k}$
    \item $\div\vec{F} = \pdv{F_x}{x} + \pdv{F_y}{y} + \pdv{F_z}{z}$
    \item $\curl\vec{F} = \begin{vmatrix} \hat{i} & \hat{j} & \hat{k} \\ \pdv{}{x} & \pdv{}{y} & \pdv{}{z} \\ F_x & F_y & F_z \end{vmatrix}$
    \item $\laplacian\phi = \pdv[2]{\phi}{x} + \pdv[2]{\phi}{y} + \pdv[2]{\phi}{z}$
    \item $\curl(\grad\phi) = \vec{0}$
    \item $\div(\curl\vec{F}) = 0$
    \item $\div(\grad\phi) = \laplacian\phi$
\end{itemize}
\end{multicols}

\subsection{Complex Analysis Formulas}
\begin{multicols}{2}
\begin{itemize}
    \item $e^{i\theta} = \cos\theta + i\sin\theta$
    \item Cauchy-Riemann: $\pdv{u}{x} = \pdv{v}{y}$, $\pdv{u}{y} = -\pdv{v}{x}$
    \item $\oint_C f(z) \, dz = 0$ (Cauchy's theorem)
    \item $f(z_0) = \frac{1}{2\pi i}\oint_C \frac{f(z)}{z-z_0} \, dz$
    \item $\text{Res}(f,z_0) = \frac{1}{2\pi i}\oint_C f(z) \, dz$
    \item $\oint_C f(z) \, dz = 2\pi i \sum \text{Res}(f,z_k)$
\end{itemize}
\end{multicols}

\subsection{Differential Equations Formulas}
\begin{multicols}{2}
\begin{itemize}
    \item Linear ODE: $\dv{y}{x} + P(x)y = Q(x)$
    \item Integrating factor: $\mu = e^{\int P \, dx}$
    \item Characteristic eq: $ar^2 + br + c = 0$
    \item Bessel: $J_\nu(x) = \sum_{m=0}^{\infty} \frac{(-1)^m}{m!\Gamma(m+\nu+1)}\left(\frac{x}{2}\right)^{2m+\nu}$
    \item Legendre: $P_n(x) = \frac{1}{2^n n!}\dv[n]{}{x}[(x^2-1)^n]$
    \item Hermite: $H_n(x) = (-1)^n e^{x^2}\dv[n]{}{x}(e^{-x^2})$
\end{itemize}
\end{multicols}

\subsection{Linear Algebra Formulas}
\begin{multicols}{2}
\begin{itemize}
    \item $A\vec{v} = \lambda\vec{v}$ (eigenvalue problem)
    \item $\det(A - \lambda I) = 0$ (characteristic equation)
    \item $A = PDP^{-1}$ (diagonalization)
    \item $A = UDU^\dagger$ (Hermitian diagonalization)
    \item Trace = sum of eigenvalues
    \item Determinant = product of eigenvalues
\end{itemize}
\end{multicols}

\newpage

\section{Study Strategy for GATE 2026}
\begin{importantbox}
\textbf{Mastery Plan:}
\begin{enumerate}
    \item \textbf{Week 1-2:} Vector Calculus - Master all operators, theorems, and coordinate systems
    \item \textbf{Week 3-4:} Complex Analysis - Focus on Cauchy theorem, residues, and contour integration
    \item \textbf{Week 5-6:} Differential Equations - Practice Frobenius method, master special functions
    \item \textbf{Week 7-8:} Linear Algebra - Deep dive into eigenvalues, diagonalization, Hermitian matrices
    \item \textbf{Week 9+:} Integration and problem-solving - Combine all topics, solve past papers
\end{enumerate}
\end{importantbox}

\subsection{Key Points to Remember}
\begin{itemize}
    \item Mathematical Physics is the foundation - weak here means struggling everywhere
    \item Practice problems daily - understanding comes from doing
    \item Focus on physical interpretations, not just calculations
    \item Master coordinate transformations (Cartesian, cylindrical, spherical)
    \item Understand when to use which method (Frobenius, residue theorem, etc.)
    \item Memorize key formulas but understand their derivations
    \item Work through examples step-by-step
    \item Solve previous GATE papers focusing on Mathematical Physics
\end{itemize}

\vspace{2cm}
\begin{center}
\textbf{Good luck with your GATE 2026 preparation!}\\
\textit{Master Mathematical Physics, and you'll master physics.}
\end{center}

\end{document}
